\chapter[Les attentes des usagers]{Les attentes des usagers finaux: étudiants, cherhceurs et professionnels du droit}

\lettrine{L}a refonte d'une bibliothèque numérique ne peut se faire sans confronter l'outil que l'on souhaite avec les usages et les pratiques réelles des publics qu'elle est censée servir. Si l'ingénierie documentaire du Web de données tend à modéliser l'inforamtion pour les machines, la réussite de la nouvelle interface dépend d'une transition méthodologique, passant d'une approche descendante - où l'utilisateur doit s'adapter aux contraintes techniques du catalogue - vers une approche ascendante, orientée vers l'expérience des utilisateurs (UX). Cette méthode ascendante, éprouvée à la BIU Cujas à travers la réalisation d'une enquête auprès de leurs usagers, complétée par quatre entretiens qualitatifs approfondis auprès de chercheurs, a permis d'identifier précisément leurs attentes envers la nouvelle bibliothèque numérique.

\section[Diversité des publics]{Vers une bibliothèque numérique orientée usagers : s'adapter à la diversité des publics de Cujas}

La refonte de CujasNum ne saurait se limiter à une simple mise à niveau sur le plan technique. En plaçant l'usager au centre de son écosystème, cela représente un impératif stratégique important. La réussite de la future plateforme repose sur le passage d'un modèle sous forme de \enquote{répertoire de PDF} à un modèle de travail capable de répondre aux exigences techniques et méthodologiques permettant à un large public de l'utiliser, allant de l'étudiant de master au chercheur émérite. L'enjeu est de transformer l'infrastructure statique en un service dynamique soutenant la production scientifique.

L'analyse des entretiens et des retours d'enquête ont mis en évidence différentes typologies de recherche qui imposent des fonctionnalités différencies : l'enseignant-chercheur expert, le doctorant, l'étudiant et le lecteur éclairé. Ces quatre profils d'usager, utilisant la bibliothèque numérique, ont des méthodes de recherche propres. L'enseignant-chercheur est spécialiste d'une question ou d'une thématique, pratique une recherche de \enquote{fond} sur des sources rares. Ses habitudes de recherche incluent la sérendipité, l'utilisation de l'OCR comme d'un outil permettant de consulter le texte d'une façon plus agréable, et télécharge massivement des PDF pour alimenter ses prores dossiers. Le doctorant a des pratiques similaires à celui de l'enseignant-chercheur, mais exige davantage de précision sur la bibliothèque numérique utilisée. Il utilise le numérique pour la préparation pédagogique, nécessitnat des outils de capture d'écrans performant pour illustrer ses supports de cours. Pour l'étudiant, la bibliothèque numérique doit être efficace, passant par l'export bibliographique automatisé (Zotero), la rapidité d'accès, des explications précises et des guides d'aide à la prise en main de l'interface et de son contenu. Pour le lecteur éclairé, celui-ci va être attiré par les expositions virtuelles, ou une autre forme de valorisation sur des thématiques d'actualité sur lesquelles ils tomberaient par hasard pour en apprendre davantage.

Le numérique répond aujourd'hui à une tripe nécessité: logistique, scientifique et patrimoniale. Pour les chercheurs, la nouvelle bibliothèque numérique doit leur permettre d'avoir accès au contenu de la bibliothèque mais également provenant d'autres institutions françaises ou étrangères. Pour les institutions patrimoniales, l'usage du numérique représente un important levier de conservation. Face à la dégradation physique des collections papiées par le temps ou les intempéries, la version numérique offre une pérennité et une \enquote{propreté} que le document physique ne garantit plus. 
 
Actuellement, la bibliothèque numérique ne permet pas une adaptation optimale aux pratiques de la recherche, laissant présumé d'une compréhension par son public alors que ce n'est pas toujours le cas. Comme l'a expliqué Jérôme Dinet, en déconstruisant le mythe de l'autonomie des étudiants. Face aux développement des nouvelles technologies, il démontre que \enquote{l'aisance technique des nouvelles générations face aux écrans ne garantit pas leur capacité à chercher, filtrer et traiter de l'information académique}. Cette fracture d'usage entraîne des experts et des publics non-spécialiste légitime l'approche d'adapter la bibliothèque numérique en concevant des parcours différenciés et adaptés à chaque niveau de compétence. 
 
Par exemple, actuellement sur CujasNum, les retours des usagers face aux contenus proposés est un mélange d'incompréhension concernant la constitution des collections. En effet, sur le portail de CujasNum, on observe un \enquote{conflit de classification} entre des logiques thématiques et chronologiques. Il s'agit d'un flagrant contraste avec la promesse faite aux lecteurs lorsqu'il se trouve sur la page d'accueil, notamment avec les titres précis et les vignettes laissent entendre d'une forme de précision. Cette précision laisse entendre pour le lecteur qu'il trouvera un contenu spécialisé en consultant l'une des collections. Par exemple, dans la collection intitulée \enquote{Droit romain}, sans médiation éditoriale entraîne une confusion, à l'instar de la présence non-justifié pour les usagers d'auteurs du XVIIIe siècle, traitant du droit romain. Le retour fait notamment par les doctorants expliquent comme créant une collection \enquote{Droit romain}, ce n'est pas à cela qui s'attend mais plutôt à l'ensemble des lois, des règles et des principes juridiques régissant la Rome antique. Le manque d'explication derrière les choix faits par les experts métiers entraînent une confusion d'interprétation sur la logique suivie, engendrant une \enquote{peur de rater des sources}. Loin d'être confortable, c'est une des attentes majeures pour la nouvelle bibliothèque numérique de la part des usagers. Une autre de leurs attentes est de pouvoir retrouver le confort du livre papier qui est recréer numériquement. 
\section[Les attentes fonctionnelles]{Les attentes fonctionnelles : recréer numériquement le confort du livre papier}
Dans une époque où l'urgence des moteurs imposent de produire plus en détrimant de la qualité, le confort de lecture ne doit pas être perçu comme un luxe cosmétique, mais comme un socle de la \enquote{Slow Science}. Les attentes des usagers sont donc de pouvoir recréer le confort papier du livre en le développant numérique à travers les tables des matières cliquables et le développement de l'OCR. Actuellement, l'un des avantages de CujasNum pour les usagers qui en sont \enquote{grandement satisfait est la présence de la table des matières cliquables, implémentées dans leur visionneuse JPEG mais aussi dans le PDF}. La mise en place de ce service permet aux utilisateurs de mieux appréhender le défilement fastidieux des corpus volumineux. La table des matières permet comme l'expliquait un des enseignants-chercheurs d'aller directement à la section souhaitée, complété par l'OCR. Cela favorise la recherche numérique. L'OCR a été massivement adopté par les institutions patrimoniales permettant aux chercheurs, aux étudiants de faire de la recherche plein-texte sur les images numérisées. C'est un outil très précieux et utile encrer dans les pratiques de recherche d'aujourd'hui. L'absence d'océrisation condamne le portail documentaire ou le mauvais référencement de celui-ci empêche celui-ci de pouvoir être exploité comme il le devrait. 

Les chercheurs demandent s'ils seraient possible de pouvoir mettre en place pour la future bibliothèque numérique de l'OCR qui soit capable aussi bien de faire des recherches plein-texte que de traiter des occurences précises en latin également avec le surlignage de termes lorsqu'ils sont trouvés dans la liste de résultat, accompagnée d'un lien cliquable pour se rendre directement au passage identifié, ou dans la visionneuse. L'adoption massive du standard IIIF est devenu le levier de la transformation numérique. La visionnneuse doit permettre de réaliser \enquote{un zoom de haute résolution, la comparaison des documents et des outils de captures dynamiques}. Par exemple, l'ajout des outils de captures d'écran sur la visionneuse permettrait à l'un des doctorants interviewés de mieux préparer ces illustrations pédagogiques qu'il utilise sur ces powerpoints. L'ajout de visionneuse IIIF, comme \textit{Mirador} ou \textit{Universal Viewer} a été adopté par de nombreuses institutions patrimoniales au cours de ces dernières décennies permettant à ceux qui l'utilise de pouvoir récupérer le manifeste IIIF de le consulter sur Mirador en le comparant à une autre édition présent dans une autre bibliothèque. L'avantage est que la personne n'a pas besoin de se déplacer pour cela dans l'institution. 

La préservation de la sérendipité est importante pour les usagers. Inspiré par la bibliothèque de Munich, la nouvelle bibliothèque numérique de Cujas doit proposer une interface relationnelle. La \enquote{recherche bucolique} pronée par l'un des enseignants-chercheurs n'est pas une errance, mais une méthodologie historique. Le système doit suggérer des \enquote{chemins de traverses} à travers des liens vers des éditions similaires, des auteurs cités ou des documents connexes, permettant une exploration par cercles concentriques plutôt qu'une simple réponse à une requête unique. La contrainte de la recherche n'est pas quelque chose d'appréciée par les chercheurs, c'est le risque lorsqu'il fait une requête et que celle-ci retourne trop de résultats ou alors très peu. Si la vision des collections des experts métiers s'impose de trop dans la méthode de recherche concernant la vision des collections qu'elle a pour les utilisateurs \enquote{les empêchant ainsi de faire leur propre découverte}. Cela ne lui laisse alors ni l'opportunité de s'approprier les collections par plusieurs tentatives de requêtes ni de faire des découvertes inattendues. Il est important de trouver un équilibre entre \enquote{la recherche ciblée par un but précis (trouvabilité) à la recherche exploratoire par sérendipité (découvrabilité)}, comme l'a expliqué Elina Leblanc dans sa thèse. 

Un autre garant de la sérendipité numérique est le \enquote{browsing}(le feuilletage libre), recréant l'expérience physique de la déambulation dans les rayons de la bibliothèque Cujas. Face à l'augmentation des collections numérisées, dans le monde des bibliothèques numériques, constitue \enquote{à la fois un avantage pour les lecteurs, leur offrant accès à toujours plus de documents et de connaissances. Le risque est la difficulté de trouver efficacement les ressources recherchées}. L'une des solutions est le recours au \enquote{browsing} permettant à l'utilisateur de naviguer \enquote{au milieu des collections, en empruntant des chemins qui ne lui ont pas été indiqués par la bibliothèque}. L'utilisateur peut également faire un tri dans les collections en utilisant les facettes qui lui permettent de naviguer parmi les contenus. \enquote{En choisissant lui-même les facettes, l'utilisateur crée un parcours de recherche personnalisé, qui lui permet d'investir complètement l'interface et l'encourage à poursuivre sa navigation}. L'amélioration de la recherche et du confort de lecture entraînent la nécessité de nouveaux outils permettant aux utilisateurs de sélections ce qui leur convient s'appropriant ainsi mieux le matériel. 
\section[De la consultation à l'appropriation]{De la consultation à l'appropriation : la nécessité de nouveaux outils}
Le passage d'une bibliothèque numérique de première génération, sous la forme de répertoire de PDF, à une plateforme de troisième génération, s'accompagne d'une redéfinition de la posture de l'usage en le replaçant au centre. L'usager n'est plus un \enquote{lecteur-passant} passif qui consulte un document à l'écran; il aspire à devenir un acteur actif. Il aimerait être capable d'extraire, de manipuler et de s'approprier la matière la matière textuelle et bibliographique nourrissant ces propres travaux. 

Pour guider la conception de la nouvelle interface, il convient d'adopter la théorie formulé en 2017 par Antoine Courtin. Il rappelle que: \enquote{les pratiques des chercheurs, en termes de données, se limitent bien souvent à copier des données dans un document de traitement de texte}. Face à ce constat empirique, vouloir imposers des interfaces captives ou des formats complexes et hermétiquements aux méthodes de recherche des chercheurs est inutile. La bibliothèque numérique doit proposer des services permettant l'utilisation brute des données ainsi que des services d'export s'intégrant naturellement dans les chaînes de travail personnelles des chercheurs.

Cette appropriation de la donnée par les chercheurs se heurte aujourd'hui à la volativité des sessions de recherche. Dans CujasNum, l'absence de compte utilisateur permanent condamnait le chercheur à voir son panier de documents sélectionnés disparaître dès la fermeture du navigateur. Ce mode de fonctionnement oblige le chercheur à refaire inlassablement les mêmes requêtes d'une session à l'autre afin de retrouver les documents consultés pour leurs travaux. Pour surmonter cet obstacle, la théorisation de l'espace personnel par Elina Leblanc s'avère la solution idéale pour la nouvelle bibliothèque numérique, répondant ainsi aux attentes des usagers. Leblanc préconise la  création d'un espace personnel conçu comme un véritable \enquote{carnet de lecture}. L'accès permanent à ce compte personnel ayant un historique des consultations et des recherches sauvegardées agit alors : \enquote{comme une barrière à l'oubli, mais aussi comme une économie de temps, car cela évite de recommencer ses recherches}. 

Les enquêtes qualitatives réalisées auprès des historiens du droit de Cujas confirment cette analyse. Un des professeurs explique que lorsqu'il identifie une ressource intéressante, son premier réflexe est de la télécharger afin de constituer \enquote{sa propre bibliothèque} organisée en arborescence locale sur son ordinateur. Il exprime le souhait de pouvoir transposer cette habitude de recherche à l'intérieur de la bibliothèque numérique grâce à un espace personnel persistant. De même, qu'une chercheuse insiste sur le confort d'un tel espace permettant de \enquote{mettre en favoris, lui évitant d'égarer des sources patrimoniales sur lesquelles elle doit revenir plusieurs fois au cours d'un même travail}.

Enfin, cette démarche d'autonomisation exige une pluralité de formats de sortie. Comme le souligne Elina Leblanc, proposer un éventail diversifié de formats (PDG, JPEG, exports structurés) permet de couvrir des besoins d'analyse scientifique variés dépassant le simple cadre de l'interface publique de la bibliothèque numérique de l'institution. Si le PDF conserve la préférence des chercheurs interrogés pour son caractère stable, facilement manipulable et connu, l'exportation des données bibliographiques vers des logiciels de gestion de références, comme Zotero, est devenue une attente incontournable. En plus, des formats de sortie, les pratiques de recherche des historiens révèlent le confort que représenterait la possibilité d'avoir le choix entre plusieurs options de téléchargements. Les chercheurs demandent de pouvoir faire un téléchargement partiel et ciblé. Cette demande souligne le besoin de ne part avoir à télécharger l'ensemble, lorsqu'uniquement une section ou une partie de l'ouvrage les intéresse, tout en conservant si possible la page de titre du document afin de savoir d'où provient l'extrait.

L'intégration d'un bouton d'export vers Zotero ou de la référence bibliographique intégrale s'impose comme le compagnon quotidien de l'écriture scientifique. L'affichage en clair d'une référence bibliographique normalisée \enquote{prête à être copier-coller} garantit un gain de temps précieux lors de la rédaction. Pour les utilisateurs agguerris, l'implémentation d'un module d'exportation de métadonnées riches (formats Bib Tex, CSV ou JSON), inspiré du développement mis en place dans la bibliothèque numérique \textit{Lucienne} de l'ENS permet d'automatiser l'importation massifs de corpus. En facilitant, le passage de la consultation à la possibilité de manipuler les données brutes, offre à la BIU Cujas une meilleure exploitation de ses données par les scientifiques. Il est donc nécessaire de mettre en place des accompagnements leur permettant de prendre en main ces outils.
\section[Conception et accompagnement]{Concevoir les services et accompagner les usagers dans un système en mutation}
Aujourd'hui, la mise en oeuvre d'une bibliothèque numérique ne peut se limiter uniquement à l'assemblage d'un répertoire de fichiers statiques. Pour que la nouvelle plateforme s'impose comme un véritable outil scientifique, l'établissement doit mettre en place une forte médiation documentaire. Pour cela, la bibliothèque Cujas doit transformer un stock de fichiers en un environnement de connaissances éditorialisées. Aujourd'hui, le rôle des bibliothèques ne se limite pas à stocker des fichiers isolés; il s'agit d'assembler, de contextualiser et de faire dialoguer ces ressources entre elles pour en faire émerger une plus value intellectuelle et de nouvelles connaissances. Dans leur ouvrage, Antoine Isaac, Emmanuelle Bermès et Gautier Poupeau expliquent ce phénomène à travers le concept de la \enquote{citation de contenus}. Pour les auteurs, elle consiste à \enquote{sélectionner des ressources sur le Web, de les réorganiser, de les annoter et de donner à consulter ce travail dans une interface particulière permettant de saisir le sens qui émerge de la juxtaposition de ces ressources}. Dans ce troisième âge, ce traitement intellectuel est l'un des moyens utilisés par lese institutions patrimoniales, transformant le document-objet en un espace d'apprentissage et de valorisation du travail scientifique réalisés par des spécialistes. 

Le premier jalon de cette attente réside dans la réintroduction du paratexte d'accompagnement autour des oeuvres numérisées. Les constats issus des enquêtes de terrain à la BIU Cujas sont univoques: les usagers soulignent l'importance de mettre en valeur ce qui est numérisé avec un texte descriptif. Cette attente est vigoureusement relayée par les chercheurs interrogés. L'un des professeurs insiste sur le fait qu'il est capital \enquote{de présenter les ouvrages et d'expliciter leur qualité pour susciter la curiosité scientifique de publics}, permettant aux spécialistes et au grand public de s'y intéresser notamment si on lie cette éditorialisation à des notions d'actualité.

L'une des pistes envisagées, serait la création d'une \enquote{galerie des grands juristes} sous forme d'expositions virtuelles et de parcours thématiques commentés. Laurent Pfister propose \enquote{d'associer directement la communauté de recherche}ce projet d'éditorialisation, \enquote{en confiant à des doctorants ou à des jeunes chercheurs la rédaction de courts articles scientifiques}. Ces articles expliciteraient une thématique ou un sujet de recherche où le matériel utilisé fut les documents numérisés par la bibliothèque Cujas. Cette éditorialisation partagée présente un double bénéfice : elle valorise les travaux des jeunes chercheurs tout en tissant un lien entre la bibliothèque et les laboratoires de recherche.

L'éditorialisation est possible sous différentes formes en fonctions des formats ou des protocoles utilisés par la bibliothèque, façonnant l'espace numérique. Comme le souligne Elina Leblanc, elle s'appuie sur la définition de Marcello Vitali-Rosati en disant que \enquote{les protocoles employés pour numériser et partager des contenus dans un environnement numérique contribuent à définir plusieurs manières de percevoir le patrimoine  écrit numérique}. L'éditorialisation, la création de parcours thématiques guidés et l'ajout de notices d'autorité vulgarisées sont devenus des services indispensables pour accompagner l'étudiant dans l'appropriation des sources complexes. En mettant en place ce type de contenu, cela permet de répondre à l'une des attentes d'une doctorant concernant le découragement étudiant face à la masse de documents numérisés. Elina Leblanc préconise une médiation active afin de réduire la surcharge cognitive et la désorientation des lecteurs.
\bigskip

La future bibliothèque numérique de la bibliothèque Cujas doit devenir un outil de référence pour la recherche juridique, offrant un environnement où la rigueur académique est soutenue et non entravée, par la technologie. C'est à travers les suggestions faites par les chercheurs eux-mêmes, qui permettront selon eux à la bibliothèque de gagner en rigueur scientifique. Face à ces exigences, le logiciel XTF ne pouvait plus répondre à leurs attentes. C'est pourquoi le prototype Codex, développé sous Python-Django, a été pensé pour implémenter techniquement les standards issus du Web de données, notamment via des API et le protocole IIIF.